\lesson{2}{Mar 18 2022 Fri (16:26:55)}{Write parallel and perpendicular line}{Unit 1}
\label{les_2_unit_1:write_parallel_and_perpendicular_line}

Consider the line: $y = -2x + 1$.

Let's:

\begin{itemize}
  \label{item:instructions}

  \item Find the equation of the line that is parallel to this line and passes
    through the point $(7, -5)$. 
  \item Find the equation of the line that is perpendicular to this line
    through the point $(7, -5)$.
\end{itemize}

Now, to equation $y = -2x + 1$ is written in the slope-intercept form:
$y = mx + b$. In this form, the slope $m$ is $-2$.

\begin{itemize}
  \label{item:step_by_step_guide_to_parallel_and_perpendicular}

  \item We can use the \textbf{Parallel Slope Property}. Since the given lines
    has a slope of $-2$, a line parallel to it must also have the same slope,
    which is $-2$. So, the equation of the parallel line will have the form
    $y = -2x + b$. The line passes through $(7, -5)$, so we use $x = 7$ and
    $y = -5$ to solve for $b$:

    \begin{equation}
      \begin{split}
        y   &= -2x + b \\
        -5  &= -2(7) + b \\
        -5  &= -14 + b \\
        b   &= 9
      \end{split}
    .\end{equation}

    Now, we know the equation of the parallel line, which is $y = -2x + 9$. 

    \item We use the \textbf{Perpendicular Slope Property}.
      Since the given lines has the slope $-2$, a line with the slope
      $\frac{1}{2}$ is perpendicular to it.

      So, the equation of the perpendicular line will have the form
      $y = \frac{1}{2}x + b$.

      The line passes through $(7, -5)$, so we use $x = 7$ and $y = -5$ to solve
      for $b$.

      \begin{equation}
        \begin{split}
          y   &= \frac{1}{2}x + b \\
          -5  &= \frac{1}{2}(7) + b \\
          -5  &= \frac{7}{2} + b \\
          -\frac{10}{2} &= \frac{7}{2} + b \\
          b   &= -\frac{17}{2}
        \end{split}
      .\end{equation}
\end{itemize}

Equation of the parallel line: $\quad y = -2x + 9$

Equation of the perpendicular line: $\quad y = \frac{1}{2}x - \frac{17}{2}$

\newpage
